% GENERATED FILE. Edit canonical lessons, then run npm run generate:books.
% canonical-source-hash: fnv1a64:079a9228431c0925
% canonical-lessons: PT-C443-esquadra, PT-C443-quartel, PT-C443-conservatoria, PT-C443-denunciar, PT-C443-requerer

\chapter{Serviços públicos}
\label{ch:pt-servicos-publicos}
\hlchaptermodality{443}

\begin{chapteropening}
\textbf{By the end of this chapter:} \emph{I can deal with public offices: the police station, the barracks, the registry office, reporting a crime and applying for something.}

Complete the last lesson of chapter 443: I can deal with public offices: the police station, the barracks, the registry office, reporting a crime and applying for something.
\end{chapteropening}

\section[a esquadra]{a esquadra — the police station; the squadron}

\label{lesson:PT-C443-esquadra}



\begin{quote}
Before the new one: say the Portuguese for the cloister, then the Portuguese for to wander.
\end{quote}

\subsection*{You'll want to know: a esquadra}
\textbf{a esquadra} — \textquotedblleft{}the police station; the squadron\textquotedblright{}.

\begin{quote}
Fui à esquadra porque me roubaram a carteira.
\end{quote}

\begin{quote}
\emph{I went to the police station because my wallet was stolen.}
\end{quote}

In Brazil: \emph{a delegacia}.

\begin{cousinweb}
\emph{esquadra} was borrowed from Italian \emph{squadra} 'square, team', which goes back to Vulgar Latin \emph{exquadrāre} 'to make square', on \emph{quadrus} 'square' (Lewis \& Short).
\end{cousinweb}

\begin{grammarlens}[title={What you've built}]
One more word for telling, explaining and describing at B1.
\end{grammarlens}

\subsection*{Your turn}
\begin{itemize}
\raggedright
  \item \emph{Say it:} \emph{a esquadra}
  \item \emph{Say it:} \emph{a esquadra}, once more
  \item \emph{Say it:} a sentence of your own with \emph{a esquadra}
  \item \emph{Recall:} say the Portuguese for the cloister, then the Portuguese for to wander, then say \emph{a esquadra} again
\end{itemize}

\begin{tcolorbox}[breakable,colback=teal!4,colframe=teal!35!black,title={Before you move on}]
What does a esquadra mean? (\textquotedblleft{}The police station; the squadron\textquotedblright{}.) Say it once more.
\end{tcolorbox}
\section[o quartel]{o quartel — the barracks; the station (of firefighters)}

\label{lesson:PT-C443-quartel}



\begin{quote}
Before the new one: say the Portuguese for to wander, then the Portuguese for the police station.
\end{quote}

\subsection*{You'll want to know: o quartel}
\textbf{o quartel} — \textquotedblleft{}the barracks; the station (of firefighters)\textquotedblright{}.

\begin{quote}
O quartel dos bombeiros fica ao lado da escola.
\end{quote}

\begin{quote}
\emph{The fire station is next to the school.}
\end{quote}

\emph{O quartel dos bombeiros} — the fire station.

\begin{cousinweb}
\emph{quartel} came from French \emph{quartier} 'quarter, district' (EtymDB), from Latin \emph{quartārius} 'a fourth part', on \emph{quārtus} 'fourth'.
\end{cousinweb}

\begin{grammarlens}[title={What you've built}]
One more word for telling, explaining and describing at B1.
\end{grammarlens}

\subsection*{Your turn}
\begin{itemize}
\raggedright
  \item \emph{Say it:} \emph{o quartel}
  \item \emph{Say it:} \emph{o quartel}, once more
  \item \emph{Say it:} a sentence of your own with \emph{o quartel}
  \item \emph{Recall:} say the Portuguese for to wander, then the Portuguese for the police station, then say \emph{o quartel} again
\end{itemize}

\begin{tcolorbox}[breakable,colback=teal!4,colframe=teal!35!black,title={Before you move on}]
What does o quartel mean? (\textquotedblleft{}The barracks; the station (of firefighters)\textquotedblright{}.) Say it once more.
\end{tcolorbox}
\section[a conservatória]{a conservatória — the registry office}

\label{lesson:PT-C443-conservatoria}



\begin{quote}
Before the new one: say the Portuguese for the police station, then the Portuguese for the barracks.
\end{quote}

\subsection*{You'll want to know: a conservatória}
\textbf{a conservatória} — \textquotedblleft{}the registry office\textquotedblright{}.

\begin{quote}
Casaram na conservatória da cidade.
\end{quote}

\begin{quote}
\emph{They got married at the town registry office.}
\end{quote}

Births, marriages and deaths are recorded at \emph{a conservatória do registo civil}.

\begin{cousinweb}
\emph{conservatória} is built on \emph{conservar} 'to keep', from Latin \emph{cōnservāre} 'to preserve' (EtymDB): the office that keeps the records.
\end{cousinweb}

\begin{grammarlens}[title={What you've built}]
One more word for telling, explaining and describing at B1.
\end{grammarlens}

\subsection*{Your turn}
\begin{itemize}
\raggedright
  \item \emph{Say it:} \emph{a conservatória}
  \item \emph{Say it:} \emph{a conservatória}, once more
  \item \emph{Say it:} a sentence of your own with \emph{a conservatória}
  \item \emph{Recall:} say the Portuguese for the police station, then the Portuguese for the barracks, then say \emph{a conservatória} again
\end{itemize}

\begin{tcolorbox}[breakable,colback=teal!4,colframe=teal!35!black,title={Before you move on}]
What does a conservatória mean? (\textquotedblleft{}The registry office\textquotedblright{}.) Say it once more.
\end{tcolorbox}
\section[denunciar]{denunciar — to report (a crime); to denounce}

\label{lesson:PT-C443-denunciar}



\begin{quote}
Before the new one: say the Portuguese for the barracks, then the Portuguese for the registry office.
\end{quote}

\subsection*{You'll want to know: denunciar}
\textbf{denunciar} — \textquotedblleft{}to report (a crime); to denounce\textquotedblright{}.

\begin{quote}
Denunciámos o roubo na esquadra.
\end{quote}

\begin{quote}
\emph{We reported the theft at the police station.}
\end{quote}

Also to expose someone: \emph{denunciar um colega}.

\begin{cousinweb}
\emph{denunciar} is a learned borrowing from Latin \emph{dēnūntiāre} 'to announce, to give notice' (EtymDB), on \emph{nūntius} 'messenger'.
\end{cousinweb}

\begin{grammarlens}[title={What you've built}]
One more word for telling, explaining and describing at B1.
\end{grammarlens}

\subsection*{Your turn}
\begin{itemize}
\raggedright
  \item \emph{Say it:} \emph{denunciar}
  \item \emph{Say it:} \emph{denunciar}, once more
  \item \emph{Say it:} a sentence of your own with \emph{denunciar}
  \item \emph{Recall:} say the Portuguese for the barracks, then the Portuguese for the registry office, then say \emph{denunciar} again
\end{itemize}

\begin{tcolorbox}[breakable,colback=teal!4,colframe=teal!35!black,title={Before you move on}]
What does denunciar mean? (\textquotedblleft{}To report (a crime); to denounce\textquotedblright{}.) Say it once more.
\end{tcolorbox}
\section[requerer]{requerer — to apply for; to require}

\label{lesson:PT-C443-requerer}



\begin{quote}
Before the new one: say the Portuguese for the registry office, then the Portuguese for to report.
\end{quote}

\subsection*{You'll want to know: requerer}
\textbf{requerer} — \textquotedblleft{}to apply for; to require\textquotedblright{}.

\begin{quote}
Requeri um cartão novo na câmara municipal.
\end{quote}

\begin{quote}
\emph{I applied for a new card at the town hall.}
\end{quote}

\emph{O requerimento} — the application form.

\begin{cousinweb}
\emph{requerer} is inherited from Latin \emph{requīrere} 'to seek, to ask for' (EtymDB), built on \emph{quaerere} 'to seek'.
\end{cousinweb}

\begin{grammarlens}[title={What you've built}]
One more word for telling, explaining and describing at B1.
\end{grammarlens}

\subsection*{Your turn}
\begin{itemize}
\raggedright
  \item \emph{Say it:} \emph{requerer}
  \item \emph{Say it:} \emph{requerer}, once more
  \item \emph{Say it:} a sentence of your own with \emph{requerer}
  \item \emph{Recall:} say the Portuguese for the registry office, then the Portuguese for to report, then say \emph{requerer} again
\end{itemize}

\begin{tcolorbox}[breakable,colback=teal!4,colframe=teal!35!black,title={Before you move on}]
What does requerer mean? (\textquotedblleft{}To apply for; to require\textquotedblright{}.) Say it once more.
\end{tcolorbox}
